\subsubsection*{守恒条件、碰撞与质心系}

系统某方向的合外力始终为零，该方向的总动量才在整个过程中保持不变。若仅知道某段时间的
总外冲量为零，只能直接推出这段时间起点与终点的动量相等，中间可以变化。例如物块在光滑
水平面内碰撞，水平方向无外力；竖直方向虽有支持力，但支持力可与重力抵消，不能因为存在
支持力就断言竖直动量不守恒。列式前必须明确系统边界、方向及所研究的时间段。

动能定义为 $K=mv^2/2$；它是标量，不是动量的另一个名称。一维碰撞通常还需要恢复系数
\begin{equation}
 e=\frac{\text{分离相对速度}}{\text{接近相对速度}}
 =\frac{v_2-v_1}{u_1-u_2},\qquad 0\leq e\leq1,
\end{equation}
与动量守恒共同确定碰后速度。$e=1$ 为完全弹性碰撞，系统动能也守恒；$e=0$ 时两物体具有相同
末速度，是完全非弹性碰撞。一般碰撞中动能减少量转化为形变、内能和声能，不能把动能不守恒
误说成能量不守恒。

质心速度 $\vec V_C=\vec P/M$。把每个质点速度分解成 $\vec v_i=\vec V_C+\vec v_i'$，其中
$\vec v_i'$ 是相对质心的速度。由定义
\[
 \sum_i m_i\vec v_i'=\sum_i m_i\vec v_i-M\vec V_C=0.
\]
逐项展开速度平方，系统动能为
\[
 T=\sum_i\frac12m_i(V_C^2+2\vec V_C\cdot\vec v_i'+{v_i'}^2)
 =\frac12MV_C^2+\vec V_C\cdot\sum_i m_i\vec v_i'+\sum_i\frac12m_i{v_i'}^2.
\]
中间交叉项为零，故
\begin{equation}
 T=\frac12MV_C^2+\sum_i\frac12m_i{v_i'}^2.
\end{equation}
第一项是整体平动动能，第二项是质心系中的相对运动动能。无外力时 $\vec V_C$ 不变，碰撞只会
重新分配或耗散第二项。这个分解能直接给出：在总动量一定时，所有部分共同以质心速度运动时
总动能最小；完全非弹性碰撞的动能损失因此最大。

爆炸与非弹性碰撞的能量转化方向不同，但不意味着一般非弹性碰撞能够严格逆演。
若爆炸时间很短，可忽略外冲量，总动量仍守恒，但内部储存的化学能或弹性
势能会使动能增加。解题时先在方便的惯性系中写动量关系，再用相对速度或能量条件闭合；在质心
系中各碎片总动量始终为零，几何关系通常最清楚。
